\chapter{Beyond \texorpdfstring{$\mathcal{P}$Torch}{��Torch}}
$\mathcal{P}$Torch demonstrates that projection-based learning can be competitive with gradient-based methods on shallow architectures. By building natively on PyTorch's autograd engine, it removes prior memory bottlenecks and connects cyclic projections to first-order optimization through nonlinear relaxation.

However, two major obstacles remain. The first is depth: our vanishing-target theorem shows that the learning signal attenuates with network depth, creating a challenge structurally analogous to the vanishing-gradient problem that limited early deep learning before the advent of residual connections and principled initialization schemes. The second is coverage: effective projections for many task-specific architectures have yet to be explored and implemented.

Neither barrier is likely to fall in a single paper. Gradient-based deep learning succeeded not because every early idea worked, but because enough ideas were explored for the successful ones to accumulate into a practical toolkit. Projection-based learning may require a similar period of experimentation. In that setting, $\mathcal{P}$Torch provides both a practical research platform and a competitive baseline: its projections can already improve the existing system (PJAX), while its native integration with PyTorch enables hybrid architectures and direct comparison with standard optimization pipelines.